\documentclass[12pt]{article}

% ----- packages -----
\usepackage[margin=1in]{geometry}
\usepackage{amsmath, amssymb, amsthm, mathtools}
\usepackage{graphicx}
\usepackage{booktabs}
\usepackage{xcolor}
\usepackage{hyperref}
\usepackage{caption}
\usepackage{enumitem}
\usepackage{authblk}

\definecolor{myblue}{RGB}{33,84,157}
\definecolor{instr}{RGB}{200,30,30}   % red, for instructions to be deleted

\hypersetup{colorlinks=true, linkcolor=myblue, urlcolor=myblue, citecolor=myblue}

% Convenience: a command for instructional text. Wrap any guidance to students
% in \instr{...}; the text will appear in red. DELETE all such commands (or
% replace them with your own writing) before submitting.
\newcommand{\instr}[1]{{\color{instr}\itshape #1}}

% ============================================================
\title{\bfseries Dynamics of the Hassell Population Model}
\author{Damon Lewis, Mateo Merrin}
\affil{CPSC/MATH 455 -- Chaos and Discrete Dynamical Systems \\ Gonzaga University}
\date{\today}

\begin{document}
\maketitle

% ============================================================
\begin{abstract}
\instr{Write a concise (150--250 word) summary of the topic, objectives,
methods, and key findings. State the map you studied, the parameter range
you considered, the principal techniques you used (Newton's method,
bifurcation diagrams, Lyapunov exponents, symbolic dynamics, etc.), and the
main quantitative and qualitative conclusions. The abstract should stand
alone: a reader who reads only the abstract should know what you did and what
you found. Avoid citations and detailed equations. \textbf{Delete this
instruction before submitting.}}


\end{abstract}

% ============================================================
\section{Introduction}

\instr{State the problem and the purpose of the study. Provide background
and context: where does this map (or this question) sit in the broader
literature on discrete dynamical systems? Cite the original source if the map
is named after a researcher (Hassell, Ricker, May, Feigenbaum, etc.). Briefly
preview the structure of the rest of the paper: ``In Section~2 we describe
the methods used. Section~3 presents our findings on fixed-point structure,
the period-doubling cascade, and the Lyapunov exponent. Section~4 interprets
these findings in light of standard theory. We conclude in Section~5.''
Aim for roughly 200--400 words. \textbf{Delete this instruction before submitting.}}

% ============================================================
\section{Methods}

\instr{Describe the analytical and computational tools used. State the map
explicitly with its formula, domain, and parameter range. Include the
relevant equations: definitions of fixed points, the formula for the
multiplier of a periodic cycle, the Lyapunov exponent definition, the
Schwarzian derivative formula, and so on. Describe your computational
algorithms: how you generated the bifurcation diagram (number of parameter
samples, transient length, plotting length), how you computed the Lyapunov
exponent (orbit length, starting point, transient discarded), and how you
located periodic points or super-stable parameters (Newton's method with
seeds and convergence criteria). State the software environment you used
(Python with NumPy and Matplotlib, MATLAB, Mathematica). Pseudocode or
flowcharts are welcome for non-trivial routines. Aim for roughly 300--500
words. \textbf{Delete this instruction before submitting.}}

% ============================================================
\section{Results}

\instr{Present your findings, organized so that a reader can follow the
logical progression from the geometry of the map, to fixed points, to
bifurcations, to orbit diagrams and Lyapunov spectra, to symbolic dynamics
and conjugacy. The subsections below correspond to the eleven required
analyses in the project guidelines. Adjust the headings to fit your map and
your story, but address each item somewhere in this section. Use figures,
tables, and equations liberally and refer to them in the text (e.g.
``Figure~\ref{fig:bif} shows\ldots''). Aim for roughly 1{,}000--1{,}500 words
across all subsections. \textbf{Delete this instruction before submitting.}}

% ----- 1. The map and its phase space -----
\subsection{The map and its phase space}

\instr{State the map, its domain, and the role of the bifurcation parameter.
Describe the geometric features of the graph: monotonicity, location and
character of any critical points, whether the map is unimodal, smooth or
piecewise, and any symmetries. Include a plot of the graph $y = F(x)$ at a
representative parameter value with the diagonal $y = x$ overlaid.
\textbf{Delete this instruction before submitting.}}

% ----- 2. Time series and cobweb analysis -----
\subsection{Time series and cobweb analysis}

\instr{Produce time-series plots ($x_n$ versus $n$) and graphical (cobweb)
iterations for representative parameter values drawn from the periodic
regime, near the accumulation point, and in the chaotic regime. Use these
plots to motivate the analysis that follows: ``At low parameter values we
see convergence to a fixed point; near $a = a_*$ the orbit visibly
period-doubles; at $a = a_{\text{chaos}}$ orbits appear to fill an interval.''
\textbf{Delete this instruction before submitting.}}

% ----- 3. Fixed points and stability -----
\subsection{Fixed points and their stability}

\instr{Find all fixed points analytically (or numerically when closed forms
are unavailable). Classify each as attracting, repelling, or neutral by
computing the multiplier $F'(\bar x)$ and comparing to $1$ in absolute value.
If a fixed point is neutral, determine whether it is weakly attracting or
weakly repelling by examining higher-order terms. Describe how the
fixed-point structure changes as the bifurcation parameter varies (where
fixed points are born, where they exchange stability, where they collide).
\textbf{Delete this instruction before submitting.}}

% ----- 4. Periodic points and stability of cycles -----
\subsection{Periodic points and stability of cycles}

\instr{Locate cycles of low period---at minimum, cycles of period 2, 3, and
4---and analyze their stability using the derivative of the appropriate
iterate of $F$. The multiplier of a period-$n$ cycle through points
$x_0, x_1, \dots, x_{n-1}$ is $\prod_{j=0}^{n-1} F'(x_j)$. If your map admits
a closed-form analysis of the period-2 cycle (this is rare but happens for
$Q_c$, $F_c$, and a few others), include it.
\textbf{Delete this instruction before submitting.}}

% ----- 5. Bifurcation analysis and the period-doubling cascade -----
\subsection{Bifurcation analysis and the period-doubling cascade}

\instr{Identify the bifurcations the map undergoes as the parameter varies.
Locate (analytically or numerically) at least the first four
period-doubling bifurcation values $b_1, b_2, b_3, b_4, \dots$ (more if you
can). From these, estimate the Feigenbaum constant
$\delta = \lim_{n \to \infty} (b_n - b_{n-1})/(b_{n+1} - b_n)$. Compare your
estimate to the universal value $\delta \approx 4.6692$ and comment on what
this comparison suggests about the universality class your map belongs to.
\textbf{If your map does not exhibit a period-doubling cascade}---for
example, the skewed tent map---explain why and substitute a discussion of
the appropriate alternative route to chaos.
\textbf{Delete this instruction before submitting.}}

% ----- 6. Orbit diagram -----
\subsection{Orbit (bifurcation) diagram}

\instr{Generate the orbit diagram over a parameter range that contains the
full cascade and the chaotic region. Discuss the principal features: the
period-doubling cascade, the accumulation point, the chaotic bands and their
mergings, and any periodic windows (in particular, the period-3 window if
present). Annotate the diagram with the bifurcation values from
Section~3.5 if possible.
\textbf{Delete this instruction before submitting.}}

% Example figure placeholder; replace with your actual figure.
% \begin{figure}[h]
%   \centering
%   \includegraphics[width=0.85\linewidth]{your_bif_diagram.pdf}
%   \caption{Bifurcation diagram for [your map] showing the
%   period-doubling cascade and the chaotic regime.}
%   \label{fig:bif}
% \end{figure}

% ----- 7. Lyapunov exponents -----
\subsection{Lyapunov exponents}

\instr{Compute the Lyapunov exponent $\lambda(a)$ as a function of the
bifurcation parameter $a$ and plot it alongside (or aligned with) your
orbit diagram. Discuss what the sign of $\lambda$ tells you, where $\lambda$
crosses zero, and how the Lyapunov plot corroborates---and refines---the
picture given by the orbit diagram. How well do the zeros of $\lambda(a)$
correspond to the bifurcation values you identified?
\textbf{Be careful with your language}: a positive Lyapunov exponent is a
strong indicator of sensitive dependence on initial conditions, but is not,
by itself, a definition of chaos.
\textbf{Delete this instruction before submitting.}}

% ----- 8. Histograms -----
\subsection{Histograms (frequency distributions)}

\instr{For at least two parameter values in the chaotic regime, generate
histograms of long orbits and discuss their qualitative features (peaks,
gaps, support). Address the limitations of histograms: how does bin width
affect what you see, and what can a histogram suggest but not prove about
periodicity or ergodicity?
\textbf{Delete this instruction before submitting.}}


% ----- 10. Schwarzian derivative -----
\subsection{Schwarzian derivative}

\instr{Compute the Schwarzian derivative
\[
SF(x) \;=\; \frac{F'''(x)}{F'(x)} - \frac{3}{2}\!\left(\frac{F''(x)}{F'(x)}\right)^{\!2}
\]
where $F$ is differentiable enough for it to make sense. Determine whether
$SF < 0$ on the relevant part of the domain and discuss the consequences. In
particular, if $SF < 0$, Singer's theorem bounds the number of attracting
cycles and constrains where they can live. \textbf{If your map is not smooth
enough for $SF$ to be defined}---for example, the cusp map or the skewed
tent map---discuss this.
\textbf{Delete this instruction before submitting.}}

% ----- 11. Creative element -----
\subsection{A creative element}

\instr{Pose and answer at least one question about your map that was not
explicitly addressed in class. Examples: comparing two histograms
quantitatively (Kolmogorov--Smirnov test, mean and variance, KL divergence);
investigating the density of periodic windows in the chaotic region;
numerically estimating the topological entropy; exploring the fractal
structure of the bifurcation diagram (box-counting dimension); comparing
your map to a related map from the literature (e.g.\ Ricker vs.\ Hassell,
or quadratic vs.\ logistic). The point of this section is to demonstrate
independent thinking---there is no single ``right'' question to ask.
\textbf{Delete this instruction before submitting.}}

% ============================================================
\section{Discussion}

\instr{Interpret your results. What do they tell you about the dynamics of
the map? How do your findings compare with general theory (Sharkovskii's
theorem on the ordering of periods, Singer's theorem on attracting cycles,
Devaney's definition of chaos, Li--Yorke's theorem on period~3 implying
chaos, Feigenbaum's universality)? How do they compare with other maps
studied in class? Address the limitations of your analysis: where did
double precision give out, where were the seeds for Newton's method
unreliable, where did transients confound the histograms? Suggest
directions for further work. Aim for roughly 300--500 words.
\textbf{Delete this instruction before submitting.}}

% ============================================================
\section{Conclusion and Future Work}

\instr{Summarize the main findings of the report in 150--300 words. Point
to extensions or open questions: ``A natural next step would be\ldots,'' or
``It would be interesting to investigate whether\ldots'' Avoid introducing
new results in this section---it should crystallize what you have already
established.
\textbf{Delete this instruction before submitting.}}

% ============================================================
% References
% Use either a manual bibliography (as below) or BibTeX/biber.
% A consistent style is required; the example below uses a numerical
% bracket style.
\begin{thebibliography}{99}

\bibitem{devaney}
R.~L.\ Devaney, \emph{An Introduction to Chaotic Dynamical Systems}, 3rd ed.,
CRC Press, 2022.

\bibitem{alligood}
K.~T.\ Alligood, T.~D.\ Sauer, and J.~A.\ Yorke, \emph{Chaos: An Introduction
to Dynamical Systems}, Springer, 1996.

\bibitem{peitgen}
H.-O.\ Peitgen, H.\ J\"urgens, and D.\ Saupe, \emph{Chaos and Fractals: New
Frontiers of Science}, 2nd ed., Springer, 2004.

% \instr{Add additional references for any sources you cite---original
% research papers, software documentation, textbook chapters. Use a
% consistent style throughout. \textbf{Delete this instruction before
% submitting.}}

\end{thebibliography}

% ============================================================
\appendix
\section{Code listings (optional)}

\instr{Include any code, lengthy derivations, or supplementary figures here.
The submitted code archive is the primary deliverable for reproducibility;
the appendix is for selected listings that the reader might want to consult
without leaving the document. \textbf{Delete this instruction before
submitting, or remove this appendix entirely if you have no supplementary
material.}}

\end{document}
